% Generated draft: requires mathematical and source review.
\documentclass[11pt]{article}
\usepackage[margin=25mm]{geometry}
\usepackage{fontspec}
\setmainfont{Latin Modern Roman}
\usepackage{amsmath,amssymb,mathtools,unicode-math}
\setmathfont{Latin Modern Math}
\usepackage{enumitem,needspace,xurl,hyperref}
\setlength{\parindent}{0pt}
\setlength{\parskip}{5pt}
\setlength{\emergencystretch}{3em}
\newcommand{\N}{\mathbb{N}}
\newcommand{\Z}{\mathbb{Z}}
\newcommand{\Q}{\mathbb{Q}}
\newcommand{\R}{\mathbb{R}}
\newcommand{\C}{\mathbb{C}}
\newcommand{\F}{\mathbb{F}}
\newcommand{\A}{\mathbb{A}}
\newcommand{\PP}{\mathbb{P}}
\newcommand{\E}{\mathbb{E}}
\newcommand{\eps}{\varepsilon}
\newcommand{\dd}{\,\mathrm{d}}
\begin{document}
\section*{UnsolvedMath ID: 30002347}
\textbf{Source problem code:} catalogue entry
\subsection{Definitions and mathematical statement}
Let \(\mathbb N_0=\{0,1,2,\ldots\}\). For a finitely generated group \(G\), let \(S\subseteq G\) be a finite symmetric generating set, meaning \(S=S^{-1}\) and \(\langle S\rangle=G\). Define
\[
 |g|_S=\min\{k\in\mathbb N_0:g=s_1\cdots s_k,\ s_i\in S\},\qquad
 v_{G,S}(n)=\bigl|\{g\in G:|g|_S\le n\}\bigr|,
\]
where the empty product is the identity. Thus \(v_{G,S}(0)=1\). Finite groups are included; the supplied statement imposes no infinitude assumption.

For functions \(f,h:\mathbb N_0\to(0,\infty)\), adopt the growth comparison
\[
 f\preccurlyeq h\quad\Longleftrightarrow\quad
 \exists C\in\mathbb Z_{\ge1}\ \forall n\in\mathbb N_0,
 \quad f(n)\le C h(Cn),
\]
and write \(f\simeq h\) when both \(f\preccurlyeq h\) and \(h\preccurlyeq f\). Write \(v_G\) for the resulting growth equivalence class of \(v_{G,S}\), which is independent of the finite generating set.

\textbf{Classification question.} Characterize the functions \(f:\mathbb N_0\to(0,\infty)\), up to \(\simeq\), for which there exist a finitely generated group \(G\) and a finite symmetric generating set \(S\) such that
\[
 f\simeq v_{G,S}.
\]

\textbf{Infimal growth question.} With the pointwise interpretation of the infimum, define
\[
 \underline v_G(n)=\inf\{v_{G,S}(n):S\subseteq G\text{ is a finite symmetric generating set}\},
 \qquad n\in\mathbb N_0.
\]
Determine how \(\underline v_G\) can differ, in growth, from \(v_{G,S_0}\), as \(G\) ranges over finitely generated groups and \(S_0\) is any fixed finite symmetric generating set of \(G\).

\textbf{Unresolved qualifications.} The supplied record does not define growth equivalence, the generating sets allowed in the infimum, or the order of taking that infimum and passing to growth equivalence classes. The definitions above give an explicit conventional reading using word balls and a pointwise infimum over all finite symmetric generating sets, without a cardinality bound. The phrase ``how far'' specifies no particular discrepancy invariant, inequality, or extremal target; consequently the source supplies a classification question rather than a definite quantitative conjecture.
\subsection{Short English statement}
Classify the possible word growth types of finitely generated groups. Also determine how their usual word growth can compare with the pointwise minimum over all finite symmetric generating sets. The conventions and the intended measure of discrepancy require source confirmation. This is a statement-only draft requiring human mathematical review.
\subsection{Sources}
\begin{itemize}
\item Supplied local UnsolvedMath registry record, permanent ID \(30002347\), title \emph{Growth Functions of Finitely Generated Groups}, source problem code \texttt{catalogue entry}, in the Oberwolfach Reports Open Problems collection. No individual proposer is identified. The record gives proposed year 2013.
\item Supplied workshop attribution: \emph{Geometric Structures in Group Theory (2014)}, Oberwolfach Reports, pp.\ 1629--1675, DOI \texttt{10.4171/owr/2013/28}: \url{https://doi.org/10.4171/owr/2013/28}. The supplied background describes these reports as documenting mathematical workshops and their open-problem discussions. The citation's 2014 label and the record's 2013 year require reconciliation.
\item Bartholdi--Erschler, \emph{Groups of given intermediate word growth}: \url{https://arxiv.org/abs/1110.3650}.
\item Erschler, \emph{Boundary behavior for groups of subexponential growth}: \url{https://annals.math.princeton.edu/2004/160-3/p09}.
\item Bartholdi, \emph{A Wilson group of non-uniformly exponential growth}: \url{https://arxiv.org/abs/math/0210471}.
\item Literature assessment dated \texttt{2026-08-21}: this is supplied dated curation, not independently verified. It labels the problem open and reports that growth types and rates are understood in many classes, that constructions realize broad intermediate-growth families, and that examples of non-uniform exponential growth exist; it further reports that the requested classification and the general relation between infimal and ordinary growth functions remain open. The supplied background calls this a dated literature triage derived from OWR curation and requests independent verification of time-sensitive status and cited claims.
\item The registry labels the statement \texttt{corrected\_verified} and reports a manual semantic, typographic, and LaTeX review of a meaning-preserving cleanup. These are supplied editorial claims, not independent mathematical or source verification in this draft.
\end{itemize}
\end{document}
